\section{Implementation}
\label{sec:implementation}

\sys is an end-to-end distributed serving system for LLMs with a placement algorithm module, a RESTful API frontend, an orchestration layer, and a parallel execution engine. The algorithm module, frontend, and orchestration layer are implemented with 6.5K lines of Python code. The parallel execution engine is implemented with 8.1K lines of C++/CUDA code.

The placement algorithm module implements the algorithm and the simulator mentioned in \S\ref{sec:method} which gives the placement decision for a specific model and cluster setting. The frontend supports an OpenAI API-compatible interface where clients can specify the sampling parameters like maximum output length and temperature. The orchestration layer manages the prefill and decoding instances, responsible for request dispatching, KV cache transmission, and results delivery. It utilizes NCCL~\cite{nccl} for cross-node GPU communication and asynchronous CudaMemcpy for intra-node communication, which avoids blocking the GPU computation during transmission. Each instance is powered by a parallel execution engine, which uses Ray~\cite{ray} actor to implement GPU workers that execute the LLM inference and manage the KV Cache in a distributed manner. It integrates many recent LLM optimizations like continuous batching~\cite{yu2022orca}, FlashAttention~\cite{dao2022flashattention}, PagedAttention~\cite{vllm} and supports popular open-source LLMs such as OPT~\cite{zhang2022opt} and LLaMA~\cite{touvron2023llama}.
